% handoutC.tex — HarpChords: regular chord vocabulary × 7 modes, key of C
% Compile with: xelatex handoutC.tex

\documentclass[12pt,letterpaper]{article}

\usepackage[margin=0.3in,bottom=0.35in]{geometry}
\usepackage{fontspec}
\usepackage{array}
\usepackage{booktabs}
\usepackage{xcolor}
\usepackage{colortbl}
\usepackage{parskip}

\setmainfont{DejaVu Serif}
\setmonofont{DejaVu Sans Mono}
\newfontfamily\symfont{DejaVu Sans}           % Unicode-rich font for ⚠
\newfontfamily\chordfont{TeX Gyre Pagella}[Numbers=OldStyle]  % Music-notation serif with oldstyle figures

% Chord-name macros
\newcommand{\rn}[2]{\textbf{#1}#2}      % bold Roman numeral + quality
\newcommand{\qs}[1]{\textit{s}#1}       % sus: italic s + N
\newcommand{\qq}{\textit{q}}            % quartal q (italic)
\newcommand{\qqs}{\textit{q}7}          % quartal 7: italic q + 7

\setlength{\parindent}{0pt}
\pagestyle{empty}

% ─── Neutral palette ──────────────────────────────────────────────────
\definecolor{headerbg}{HTML}{2C3E50}
\definecolor{altrow}{HTML}{F0F4F8}
\definecolor{desccol}{HTML}{2A3342}
\definecolor{flavcol}{HTML}{6B7A8F}
\definecolor{rhcol}{HTML}{1F4E79}
\definecolor{lhcol}{HTML}{7B2B2B}
\definecolor{warn}{HTML}{C0322A}

% ─── Rainbow: 7 mode column tints (pastel) ────────────────────────────
\definecolor{mode1}{HTML}{E8D5D5}
\definecolor{mode2}{HTML}{E8DDD0}
\definecolor{mode3}{HTML}{E8E5CB}
\definecolor{mode4}{HTML}{D5E8D0}
\definecolor{mode5}{HTML}{D0E0E8}
\definecolor{mode6}{HTML}{D5D5E8}
\definecolor{mode7}{HTML}{E0D5E8}

% ─── Rainbow: 7 header tints (saturated) ──────────────────────────────
\definecolor{hdr1}{HTML}{C0605A}
\definecolor{hdr2}{HTML}{C08040}
\definecolor{hdr3}{HTML}{B0A030}
\definecolor{hdr4}{HTML}{50A050}
\definecolor{hdr5}{HTML}{4080A0}
\definecolor{hdr6}{HTML}{5050A0}
\definecolor{hdr7}{HTML}{8050A0}

\color{desccol}

\begin{document}

\begin{center}
{\LARGE\bfseries\color{rhcol} HarpChords --- Regular Chord Vocabulary}\\[4pt]
{\small\itshape\color{flavcol} Diatonic chords in Roman numerals across 7 church modes --- sorted by shape}
\end{center}

\vspace{4pt}

% ─── Main chord table ────────────────────────────────────────────────
\renewcommand{\arraystretch}{1.1}
\setlength{\tabcolsep}{4pt}
\footnotesize
\begin{center}
\begin{tabular}{@{}
  >{\ttfamily\columncolor{altrow}}b{9mm}
  >{\columncolor{altrow}\ttfamily}b{17mm}
  >{\columncolor{mode1}\chordfont\centering\arraybackslash}p{18mm}
  >{\columncolor{mode2}\chordfont\centering\arraybackslash}p{18mm}
  >{\columncolor{mode3}\chordfont\centering\arraybackslash}p{18mm}
  >{\columncolor{mode4}\chordfont\centering\arraybackslash}p{18mm}
  >{\columncolor{mode5}\chordfont\centering\arraybackslash}p{22mm}
  >{\columncolor{mode6}\chordfont\centering\arraybackslash}p{18mm}
  >{\columncolor{mode7}\chordfont\centering\arraybackslash}p{18mm}
@{}}
\toprule
\cellcolor{headerbg}\color{white}\normalfont\fontsize{14}{16}\selectfont$\looparrowright$\raisebox{0.3ex}{\footnotesize\textbf{+}} &
\cellcolor{headerbg}\color{white}\normalfont Shape &
\cellcolor{hdr1}\color{white}\shortstack{\bfseries 1\\\normalfont\scriptsize Ionian} &
\cellcolor{hdr2}\color{white}\shortstack{\bfseries 2\\\normalfont\scriptsize Dorian} &
\cellcolor{hdr3}\color{white}\shortstack{\bfseries 3\\\normalfont\scriptsize Phrygian} &
\cellcolor{hdr4}\color{white}\shortstack{\bfseries 4\\\normalfont\scriptsize Lydian} &
\cellcolor{hdr5}\color{white}\shortstack{\bfseries 5\\\normalfont\scriptsize Mixolydian} &
\cellcolor{hdr6}\color{white}\shortstack{\bfseries 6\\\normalfont\scriptsize Aeolian} &
\cellcolor{hdr7}\color{white}\shortstack{\bfseries 7\,{\symfont ⚠}\\\normalfont\scriptsize Locrian} \\
\midrule
24                 & 1\,2\,5                & \rn{I}{\qs{2}}        & \rn{ii}{\qs{2}}        & \rn{iii}{\qs{2}}        & \rn{IV}{\qs{2}}        & \rn{V}{\qs{2}}        & \rn{vi}{\qs{2}}        & \rn{vii}{°\qs{2}}       \\
33                 & 1\,3\,5                & \rn{I}{}              & \rn{ii}{}              & \rn{iii}{}              & \rn{IV}{}              & \rn{V}{}              & \rn{vi}{}              & \rn{vii}{°}             \\
42                 & 1\,4\,5                & \rn{I}{\qs{4}}        & \rn{ii}{\qs{4}}        & \rn{iii}{\qs{4}}        & \rn{IV}{\qs{4}}        & \rn{V}{\qs{4}}        & \rn{vi}{\qs{4}}        & \rn{vii}{°\qs{4}}       \\
44                 & 1\,4\,7                & \rn{I}{\qq}           & \rn{ii}{\qq}           & \rn{iii}{\qq}           & \rn{IV}{\qq}           & \rn{V}{\qq}           & \rn{vi}{\qq}           & \rn{vii}{°\qq}          \\
332                & 1\,3\,5\,6             & \rn{I}{6}             & \rn{ii}{6}             & \rn{iii}{6}             & \rn{IV}{6}             & \rn{V}{6}             & \rn{vi}{6}             & \rn{vii}{°6}            \\
333                & 1\,3\,5\,7             & \rn{I}{7}             & \rn{ii}{7}             & \rn{iii}{7}             & \rn{IV}{7}             & \rn{V}{7}             & \rn{vi}{7}             & \rn{vii}{ø7}            \\
334                & 1\,3\,5\,8             & \rn{I}{+8}            & \rn{ii}{+8}            & \rn{iii}{+8}            & \rn{IV}{+8}            & \rn{V}{+8}            & \rn{vi}{+8}            & \rn{vii}{°+8}           \\
335                & 1\,3\,5\,9             & \rn{I}{9$_{7}$}       & \rn{ii}{9$_{7}$}       & \rn{iii}{9$_{7}$}       & \rn{IV}{9$_{7}$}       & \rn{V}{9$_{7}$}       & \rn{vi}{9$_{7}$}       & \rn{vii}{°9$_{7}$}      \\
337                & 1\,3\,5\,11            & \rn{I}{11$_{79}$}     & \rn{ii}{11$_{79}$}     & \rn{iii}{11$_{79}$}     & \rn{IV}{11$_{79}$}     & \rn{V}{11$_{79}$}     & \rn{vi}{11$_{79}$}     & \rn{vii}{°11$_{79}$}    \\
339                & 1\,3\,5\,13\,\dag      & \rn{I}{13$_{79b}$}    & \rn{ii}{13$_{79b}$}    & \rn{iii}{13$_{79b}$}    & \rn{IV}{13$_{79b}$}    & \rn{V}{13$_{79b}$}    & \rn{vi}{13$_{79b}$}    & \rn{vii}{°13$_{79b}$}   \\
353                & 1\,3\,7\,9             & \rn{I}{9$_{5}$}       & \rn{ii}{9$_{5}$}       & \rn{iii}{9$_{5}$}       & \rn{IV}{9$_{5}$}       & \rn{V}{9$_{5}$}       & \rn{vi}{9$_{5}$}       & \rn{vii}{ø9$_{5}$}      \\
355                & 1\,3\,7\,11            & \rn{I}{11$_{59}$}     & \rn{ii}{11$_{59}$}     & \rn{iii}{11$_{59}$}     & \rn{IV}{11$_{59}$}     & \rn{V}{11$_{59}$}     & \rn{vi}{11$_{59}$}     & \rn{vii}{ø11$_{59}$}    \\
357                & 1\,3\,7\,13\,\dag      & \rn{I}{13$_{59b}$}    & \rn{ii}{13$_{59b}$}    & \rn{iii}{13$_{59b}$}    & \rn{IV}{13$_{59b}$}    & \rn{V}{13$_{59b}$}    & \rn{vi}{13$_{59b}$}    & \rn{vii}{ø13$_{59b}$}   \\
444                & 1\,4\,7\,10            & \rn{I}{\qqs}          & \rn{ii}{\qqs}          & \rn{iii}{\qqs}          & \rn{IV}{\qqs}          & \rn{V}{\qqs}          & \rn{vi}{\qqs}          & \rn{vii}{°\qqs}         \\
\bottomrule
\end{tabular}
\end{center}
\normalsize

\vspace{4pt}

\textbf{\color{rhcol}Shape.} \ The ascending diatonic-interval pattern between fingers. \texttt{33} = start, up a 3rd, up a 3rd $\to$ \texttt{1\,3\,5}. Digit \texttt{3} = up a 3rd (two scale degrees); \texttt{5} = up a 5th (four scale degrees). Shape digit count = finger count $-$ 1. Hex digits \texttt{a}=10, \texttt{b}=11, \texttt{c}=12.

\vspace{4pt}

% ─── Legend: three side-by-side columns, one section each ─────────────
\footnotesize
\begin{minipage}[t]{0.32\textwidth}
\begin{tabular}{@{} >{\ttfamily}l l @{}}
\arrayrulecolor{headerbg}\toprule
\multicolumn{2}{@{}l}{\bfseries\color{rhcol} Quality} \\
\midrule
(none), M   & major triad \\
m, -        & minor \\
°, dim      & diminished \\
aug         & augmented $\ast$ \\
△, maj7, M7 & major 7 \\
m7, -7      & minor 7 \\
7           & dominant 7 \\
ø, ø7       & half-dim 7 \\
°7, dim7    & fully-dim 7 $\ast$ \\
m△7, -△7    & minor-major 7 $\ast$ \\
\bottomrule\arrayrulecolor{black}
\end{tabular}
\end{minipage}%
\hfill
\begin{minipage}[t]{0.32\textwidth}
\begin{tabular}{@{} >{\ttfamily}l l @{}}
\arrayrulecolor{headerbg}\toprule
\multicolumn{2}{@{}l}{\bfseries\color{rhcol} Sus / Quartal / Add} \\
\midrule
sus2, s2    & sus 2 \ (1\,2\,5) \\
sus4, s4    & sus 4 \ (1\,4\,5) \\
q           & quartal \ (1\,4\,7) \\
q7          & quartal 7 \ (1\,4\,7\,10) \\
6           & triad + 6 \\
+8          & triad + octave root \\
9$_{7}$     & triad + 9 \\
11$_{79}$   & triad + 11 \\
13$_{79b}$  & triad + 13 \\
\bottomrule\arrayrulecolor{black}
\end{tabular}
\end{minipage}%
\hfill
\begin{minipage}[t]{0.32\textwidth}
\begin{tabular}{@{} >{\ttfamily}l l @{}}
\arrayrulecolor{headerbg}\toprule
\multicolumn{2}{@{}l}{\bfseries\color{rhcol} Extensions \& Omissions} \\
\midrule
7           & diatonic 7 \ (1\,3\,5\,7) \\
9$_{5}$     & jazz 9 \ (1\,3\,7\,9) \\
11$_{59}$   & jazz 11 \ (1\,3\,7\,11) \\
13$_{59b}$  & jazz 13 \ (1\,3\,7\,13) \\
$_{3}$      & no 3 \ (power chord) \\
$_{5}$      & no 5 \ (\texttt{C$_5$} = dyad) \\
$_{7}$      & no 7 \ (= triad + \emph{n}) \\
$_{79}$     & no 7, 9 \\
$_{79b}$    & no 7, 9, 11 \ ($b{=}11$) \\
\bottomrule\arrayrulecolor{black}
\end{tabular}
\end{minipage}
\normalsize

\smallskip

{\color{flavcol} \textbf{Note.}\ \ The "\texttt{7}" row spans distinct qualities per mode: \texttt{I△ ii7 iii7 IV△ V7 vi7 viiø7} (dominant only on V). Same pattern for 9 / 11 / 13. \ \ Chromatic (not in this diatonic table, reference only): \texttt{aug}, \texttt{°7}, \texttt{m△7}, \texttt{$\flat$/$\sharp$}, \texttt{alt}. \ \ $\ast$ chromatic. \ \ \dag\ span 12, needs long reach or 2-hand voicing.}

\end{document}
